\documentclass[10pt]{article}
\usepackage[a4paper,margin=0.75in]{geometry}
\linespread{0.99}
\usepackage{parskip}
\usepackage[hidelinks]{hyperref}
\usepackage{xurl}
\usepackage{enumitem}
\pagestyle{empty}

\begin{document}

\noindent
Quang-Vinh Dang\\
British University Vietnam\\
Hung Yen, Vietnam\\
\href{mailto:vinh.dq4@buv.edu.vn}{vinh.dq4@buv.edu.vn}\\
ORCID \href{https://orcid.org/0000-0002-3877-8024}{0000-0002-3877-8024}

\vspace{0.7em}
\noindent 30 July 2026

\vspace{0.7em}
\noindent
The Editors\\
\emph{Machine Learning}\\
Springer

\vspace{0.8em}
\noindent Dear Editors,

\vspace{0.5em}
I submit for your consideration a research article entitled \textbf{``When
Conformal Fusion Cannot Be Calibrated: Attainability Limits for Trace-Based
Detection of Prompt-Injection Compromise.''}

The paper is a methodological contribution about combining conformal $p$-values,
motivated and evaluated on a security problem. Language-model agents that call
external tools ingest untrusted content and fold it into the context that drives
their next action, so an instruction hidden in a document can be followed as though
the user had issued it. Runtime monitors for this compute several trace-derived
signals and merge them into one alarm. Every paper I am aware of treats the merging
step as plumbing. The central result here is that it is the binding constraint, and
that the constraint is a sample-size law.

\vspace{0.4em}
\noindent\textbf{Main result}

\vspace{0.3em}
A conformal $p$-value computed from $n$ calibration values cannot fall below
$1/(n{+}1)$. That floor caps the evidence any calibrator in the standard family
$f_\lambda(p) = \lambda p^{\lambda-1}$ can express, at exactly $e^{u-1}/u$ with
$u = \log(n{+}1)$, attained at $\lambda = 1/u$. Averaging calibrated e-values is
the only merging rule that remains valid when the signals are arbitrarily
dependent, and an average never exceeds its largest term. The level-$\alpha$
dependence-robust test therefore has an \emph{empty} rejection region---power
exactly zero, on every conceivable data set---unless $e^{u-1}/u \ge 1/\alpha$. At
$\alpha = 0.05$ that first holds at $n = 312$ calibration examples. A companion
proposition gives the corresponding turn-on for the Bonferroni-style rules,
$n \ge K/\alpha - 1$, which is $39$ for two signals.

Both statements are elementary once posed, and I have not found either in the
conformal-prediction or e-value literature, which studies merging functions under
the assumption that the underlying $p$-values are continuous. Discreteness is
unavoidable for split-conformal $p$-values, and the consequence is a hard
sample-size requirement rather than a loss of efficiency.

\vspace{0.4em}
\noindent\textbf{Supporting empirical work}

\vspace{0.3em}
On $90$ genuinely compromised trajectories collected from two commercial agents,
eight combination rules produce pooled AUROC within $0.010$ of one another while
their a-priori true-positive rates span $0.000$ to $0.867$. Discrimination is
nearly rule-independent; the ability to honour a stated false-alarm rate is almost
entirely determined by the rule, which is an axis a literature reporting AUROC
alone cannot see. The three dependence-robust rules behave exactly as the
propositions require---the e-average, mean-$p$ and the union bound are inert, and
Bonferroni is the only one that fires, with power equal to the fraction of
positives carrying one signal at the $p$-value floor, $0.416$, to three decimals.
That is an identity rather than a coincidence: the two signals are at the floor
together in $0.000$ of cases. Injection evidence in these traces is therefore
dense in the bulk, where additive fusion ranks best, and sparse in the tail, where
only a minimum-based rule can be calibrated, and those two facts select opposite
rules.

The argument also explains, from one cause, a negative result the project had
recorded without explanation: an anytime-valid e-process detector was outperformed
by a hand-picked benign quantile. Time-uniform guarantees and dependence-robust
merging both purchase safety with a multiplicative budget, and both need more
calibration evidence than corpora of this size can express.

\vspace{0.4em}
\noindent\textbf{Why \emph{Machine Learning}}

\vspace{0.3em}
The contribution is statistical methodology---conformal $p$-values, e-values and
their combination under dependence---and its audience is readers who care about
calibration and validity rather than about one application domain. The security
setting supplies the motivation, a real corpus and a concrete consequence for
people building such corpora, but the propositions apply to any conformal monitor
that fuses several scores, which is now a common construction in machine-learning
systems monitoring. The paper also reports its negative results in full: five
components implemented and rejected on the evidence, a proof that the intuitive
normalisation of subtracting the largest benign training value caps AUROC at one
half, and a label-polarity trap in a standard evaluation harness that inverted my
own conclusions for a period.

Every number, table and figure is regenerated by released scripts from released
data, without any model API access, because the model-judge scores are cached and
committed. A verification script re-derives each reported value and fails on any
mismatch.

\vspace{0.4em}
\noindent\textbf{Declarations}

\vspace{0.3em}
The manuscript is original, has not been published previously, and is not under
consideration elsewhere. As sole author I am responsible for all aspects of the
work and I approve it for submission. There are no competing interests and no
external funding was received. The work involves no human participants, human data
or animals. All code and data are public:\\
\url{https://github.com/vinhqdang/SENTRY-Anytime-Valid-E-Process-Monitoring-of-LLM-Agent-Action-Streams}

\vspace{0.6em}
\noindent Thank you for your consideration.

\vspace{0.5em}
\noindent Sincerely,

\vspace{0.4em}
\noindent Quang-Vinh Dang

\end{document}
